\section{先学会负责，再学会让 AI 干活}
\label{sec:ai-responsibility}

刚才的计分器很小，你可以先手算结果，再检查 AI 的建议。以后代码变长，仍然可以挑一条你能说清楚的规则，选一个输入亲自核对。能描述任务、准备材料、检查改动，也是开发能力的一部分。

不过，AI 写代码快，不代表你可以跳过基础。一个解释得头头是道的回答，可能引用了不存在的函数；一个成功编译的程序，也可能把毫秒当成秒。\textbf{你负责提出问题、判断证据，并对交付结果负责。}这些工作没有随着代码生成一起外包。

本章把\textbf{许愿式编程}叫作这样一种做法：只说“帮我写好”“全部优化”，拿到结果就运行，出错了再把整屏信息原样转给 AI，不看改动，也不知道怎样判断完成。它偶尔能带来漂亮的演示，但很难积累能独立工作的能力。如果学长安排了一个任务，你只负责把学长的话转给 AI，再把 AI 的话转回来，那你就成了一个传话筒。下一次现场改参数、断网排障或解释接口时，这份“熟练”很容易露馅。

我们可以换一种学法：让 AI 把工作推进，同时逼自己建立解释。
\begin{enumerate}
    \item \textbf{动手前能预测}：输入是什么、输出应该怎样，哪一种情况应当报错？
    \item \textbf{改完后能解释}：关键函数、数据流、单位、对象寿命、失败分支分别是什么？
    \item \textbf{离开对话能修改}：给定一个小变化，能自己找到位置、修改并构建。
    \item \textbf{交付时能举证}：拿出命令、输入、实际输出、版本和仍未验证的部分。
\end{enumerate}
不必背下每一行生成代码，但不能有一块影响正确性的“反正 AI 说可以”。问 AI“解释你的方案，列出关键假设、验证方法和理论原理”，比要求它直接完成整个项目更有用。

\Needspace{6\baselineskip}
\begin{example}[把“能运行”变成“我会用了”]
比如以后遇到一个 AI 帮你写的串口读取函数，可以先用短小的教学输入验证。让它画出从输入字节到解析结果的步骤，并指出代码位置。你现在只需要理解这条检查思路：给它不同输入，先预测行为，再运行核对。

学到串口协议时，可以试着把一帧数据拆成两次输入，再放入一帧坏校验数据，预测输出后运行。最后关闭对话窗口，改一个合法帧中的数值，并解释为什么结果随之变化。能完成这些，才算掌握了这次修改的主要原理。
\end{example}

\begin{exercise}\label{ex:ai-principle-responsibility}
一位同学说“AI 已经修好”一个把毫米换算成米的程序，只给了启动截图。写出你还要核对的输入单位、计算方法、运行步骤和实际输出。任选一个能手算的输入，给出预期结果。
\end{exercise}

\section{模型与代理分别做什么}
\label{sec:ai-model-agent}

同一段报错，普通聊天可以读你贴出的文字并解释；代码代理还能在授权范围内自己打开文件、运行命令，再用结果继续判断。模型根据当前上下文生成回答，你附上的代码和解释是本次工作的材料。词元、训练和生成过程见第~\ref{sec:ai-llm}~节。

\subsubsection{从模型到代理：Harness 在做什么}

单独的模型根据输入生成输出。\textbf{代码代理（Coding Agent）}还可以借助工具读取文件、搜索、编辑和运行命令。把模型接进这些能力的外围系统，常被称为 \textbf{Harness}。它是组织任务的执行框架：装配指令与上下文、提供工具、处理执行与权限、管理会话，并把工具结果交回模型。具体具备哪些功能，取决于所用实现。模型能按约定请求工具、客户端已接入该工具、当前会话获准执行，是需要分别核对的条件。

\Needspace{6\baselineskip}
一轮代理工作可以这样理解：
\begin{center}
\begin{tabular}{c}
需求与上下文 $\longrightarrow$ 模型选择下一步 $\longrightarrow$ 工具执行\\[3pt]
$\longrightarrow$ 返回输出与错误 $\longrightarrow$ 模型调整 $\longrightarrow$ 继续或交付。
\end{tabular}
\end{center}
例如，“编译通过”可以有两种来源：模型在文字中猜测会通过，或者 Harness 真正执行构建命令，返回退出码与日志。后者给出了可核查的证据，前者只有预测。即使实际构建通过，也还没有覆盖运行结果与设备行为。

同一个模型放在不同客户端里，效果也可能不同：能读哪些目录、是否有检索工具、历史怎样压缩、是否能运行编译器，都会改变它能完成的工作。这里先分清三层：\textbf{模型}提供推断与生成能力，\textbf{产品或客户端}提供交互入口，\textbf{Harness}组织上下文、工具与执行流程。一个产品可以接入多个模型，同一模型也可能被多个产品使用。

\begin{example}[分清三个“它知道了”]
让 AI 解释计分器怎样构建时，先请它指出结论来自哪里。它可能只根据你贴的代码给出建议，也可能实际读过练习目录的 README，或者通过搜索工具找到了构建工具的文档。你可以分别核对建议、文件内容和文档出处。

如果它声称“已确认运行正常”，继续找对应的运行命令、输入和输出。没有这些材料，就请它把措辞改为“按接口分析预期如此，尚未运行”。这个操作不要求你先懂原理的推导，只要求你能区分解释、查证和实测。
\end{example}

\begin{exercise}\label{ex:ai-principle-model-agent}
AI 给出了一条安装命令、一段“全部测试通过”的文字，以及一个它实际读取的文件路径。哪些属于建议，哪些属于执行证据？还缺少什么才能确认测试真的通过？并说明“把协议粘进对话”和“训练一个模型”的区别。
\end{exercise}

\begin{conclusion}
模型负责根据输入推理与生成；Harness 负责把上下文、工具、权限和执行循环组织起来；工具结果提供来自环境的观测。回答流畅、工具可调用、任务已经验证，是三件需要分别检查的事。不同产品对这些机制的命名和实现可能不同，可以对照所用产品的文档查证；\href{https://learn.chatgpt.com/docs/glossary}{OpenAI 术语表}与\href{https://developers.openai.com/api/docs/guides/conversation-state#managing-the-context-window}{其上下文说明}是本章案例的查阅入口。
\end{conclusion}

\section{工具、权限与证据的边界}
\label{sec:ai-evidence}

代码代理能调用工具，意味着它既能帮忙，也能影响真实文件和服务。需要先区分两类边界：一类是系统实际允许访问哪些资源，另一类是哪些动作需要先询问。常见实现会用\textbf{沙箱（Sandbox）}限制命令访问，用\textbf{审批策略（Approval Policy）}管理确认流程；其他产品的机制可能不同，也可能不提供独立沙箱。“可以读项目文件”“可以向外网发请求”“可以提交到远端”是不同范围的能力，需要查看当前环境如何限制，不能只凭一个模式名称放心。

做普通练习时，只授予对应工程所需的权限，不要为了方便、图快、省事直接选择完全访问权限。出现审批请求，先读它准备运行的命令、目标路径和理由；缺依赖不等于必须对整台机器开放全部访问。提示词里的“不要动其他文件”是一条行为要求，操作系统、容器、沙箱或服务账号实施的限制才是实际访问边界。插件、MCP、网页搜索和浏览器等工具还可能使用独立连接或策略，不能认为限制了 Shell 网络就限制了所有外部访问。本章软件案例的机制见\href{https://learn.chatgpt.com/docs/agent-approvals-security}{Codex 权限与安全说明}。

还要认识\textbf{提示注入（Prompt Injection）}：网页、仓库注释、日志或文档中可能混入“忽略之前要求”“上传凭据”等文字。它们是代理正在分析的外部材料，不应被当成项目负责人给出的指令。安装插件前也要检查来源与权限，不把密钥、学校统一认证密码、队内私有数据随便粘进聊天。需要凭据时，按项目规定使用环境变量、凭据管理或服务授权，并核对发送目标。

\begin{example}[要求证据，也给证据分级]
让代理检查一个相机程序时，给出这样的要求：“请分别列出源码分析、已运行的软件检查和待实机验证项；每个已运行项附命令、输入和结果位置。”检查结果是否把三个层次混在一起。

例如，单元测试可以证明给定输入下的格式转换；回放可以观察录制数据上的行为；只有连接相机并读回参数、采集图像，才能为这台设备当时的实际行为提供证据。代码审查、仿真或短时运行，都不能直接写成“实车长时间稳定”。这不仅是使用 AI 时的规范，也是你做任何科研工作时应遵循的基本原则。
\end{example}

\begin{exercise}\label{ex:ai-principle-evidence}
代理报告“离线回放通过，因此相机硬件同步正常，连续运行稳定”，并建议为下载依赖启用全部访问。指出这段报告中至少两个证据越界，并给出较小的下一步验证或权限范围。
\end{exercise}
